\documentclass{presentation}

\title{Συναρτήσεις} 
\subtitle{Μονοτονία, Ακρότατα, Συμμετρία}
\author[Λόλας]{Κωνσταντίνος Λόλας }


\begin{document}

\begin{frame}
  \titlepage
\end{frame}

\begin{frame}{Ορισμοί}
  \begin{block}{Γνήσια Αύξουσα}
    Μια συνάρτηση $f:A\to\mathbb{R}$ λέγεται \emph{γνήσια αύξουσα} σε ένα διάστημα $Δ$ αν για κάθε $x_1, x_2 \in Δ$ με $x_1 < x_2$ ισχύει $f(x_1) < f(x_2)$.
  \end{block}
\end{frame}

\begin{frame}{Πώς παρατηρούμε τη μονοτονία?}

\end{frame}

\end{document}